%% Clickable relative references, and the reason there is machinery behind
%% them: a relative reference must never be a link to nowhere.
%%
%% \ref has always been a link and \Last has not, though the two spell the
%% same reference; this case pins the half that was missing.  What it is
%% really about, though, is the failure mode that comes with it.  A
%% relative reference names its target by ARITHMETIC, so it can name a
%% number no example carries -- \LLast before example 2, \Next past the
%% last one -- and a link to a destination that does not exist is not an
%% error: the backend substitutes a whole-page destination, so the click
%% lands somewhere plausible and wrong.  The engines do not agree that
%% this deserves a word (pdftex and luatex warn, in different phrasings;
%% xdvipdfmx says nothing), which is why linguexx checks it itself.
%%
%% Every number below is therefore chosen for whether its target exists:
%% RLBEFORE holds two references that resolve to no example and one that
%% does, RLPAST points past the end, and the footnote holds one of each on
%% the footnote series.  The dangling ones must reach the page as the
%% numbers they always were, must produce NO link, and must be named in
%% one warning at the end of the run.
\input{_preamble}
\usepackage{hyperref}
\begin{document}
%% -1 and 0 are examples that do not exist; 1 is the one below.
RLBEFORE \LLast\ \Last\ \Next

\ex.\label{rl:one} RLEXONE first example.
\a. RLSUBA sub a.
\b. RLSUBB sub b.
\z.

%% \Next[b] prints the letter but aims at the example: the letter anchors
%% are built from \alph and the printed letter from \Exalph, so a part
%% link would dangle in any document that renumbers its sub-examples.
RLNEXT \Next\ RLLAST \Last\ RLNNEXT \NNext\ RLPART \Next[b]
%% The p-twin forwards to \Next and must arrive at the same link: it
%% suppresses the parentheses, not the reference.
RLPTWIN \pNext

\ex. RLEXTWO second example.

%% Inside a footnote the references ride the footnote series -- except
%% \TextNext, which names the main-series example that follows the
%% footnote and must link to it and not to a footnote example.
RLFN host.\footnote{RLTEXTNEXT \TextNext
\ex. RLFNONE footnote example one.

\ex. RLFNTWO footnote example two.

RLFNLAST \Last\ RLFNNEXT \Next} done.

\ex. RLEXTHREE third example.

RLPAST \Next\ points past the last example.
\end{document}
